\documentclass[10pt,a4paper]{amsart}
\usepackage[margin=19mm]{geometry}
\usepackage{amsmath,amssymb,mathtools}
\usepackage[scr=rsfs,cal=euler]{mathalfa}
\usepackage{xfrac}
\usepackage[unicode,hidelinks,bookmarks=false]{hyperref}
\newcommand{\ie}{i.e.\ }
\newcommand{\cf}{cf.\ }
\newcommand{\resp}{respectively\ }
\newcommand{\Subsection}{\subsection}
\providecommand{\nobreakdash}{\nobreak\hbox{-}}
\setlength{\emergencystretch}{2em}
\multlinegap0pt
\usepackage{bm}
\usepackage{bbm}
\usepackage{dsfont}
\usepackage{xspace}
\usepackage{cancel}
\usepackage{mathtools}
\usepackage{amsthm}
\makeatletter
\@ifundefined{theo}{\newtheorem{theo}{Theo}}{}
\@ifundefined{lemm}{\newtheorem{lemm}[theo]{Lemma}}{}
\makeatother
\def\Bil{{\rm Bil}}
\def\Inv{{\rm Inv}}
\def\bC{\mathbb{C}}
\def\bZ{\mathbb{Z}}
\def\dg{\dot{\mathfrak{g}}}
\def\fL{\mathfrak{L}}
\def\fg{\mathfrak{g}}
\def\fl{\mathfrak{l}}
\title[$(d,\sigma)$-twisted Affine-Virasoro superalgebras]{$(d,\sigma)$-twisted Affine-Virasoro superalgebras}
\author{Rencai L\"u, Xizhou You, Kaiming Zhao}
\date{Source: arXiv:2507.00349v1 (CC BY 4.0)}
\begin{document}
\maketitle

\noindent\emph{Source and licence.} $(d,\sigma)$-twisted Affine-Virasoro superalgebras. Version arXiv:2507.00349v1 (CC BY 4.0), distributed under the Creative Commons Attribution 4.0 International licence, as recorded in the arXiv licence metadata for this exact version and shown on the abstract page https://arxiv.org/abs/2507.00349v1. Full rights notice: licenses/arxiv-2507.00349-CC-BY-4.0.txt.\par\smallskip

\noindent\textit{Editorial scope.} Counted unit: the Lemm whose source label is \texttt{lemm-10} in the article (source0.tex lines 509--509, source label \texttt{lemm-10}), delivered together with the article's own complete original proof of it (source0.tex lines 511--563, one \texttt{proof} environment of 53 lines). No mathematics is written, repaired or re-derived: statement and proof are the article's own text line for line. Required context retained uncounted: lemm-6 (lines 325--333); g-g-1 (lines 372--372); o-m (lines 384--386); lhd (lines 388--390); Eq2.13 (lines 426--426). Every definition, display and label of those lines is retained unchanged. Omitted from this package and not redistributed: 1--324, 334--371, 373--383, 387--387, 391--425, 427--508, 510--510, 564--1333. Nothing inside a retained excerpt is omitted or altered: every retained body is delivered exactly as the article's lines are written, so no omission receipt of any kind accompanies this package. Citations: the retained excerpts carry no citation command at all, so no bibliography entry of the article is transcribed and no reference is redistributed. Reference adaptation: none was needed and none was made; every reference inside the delivered unit resolves inside it, so no unresolved reference or citation is produced. Printed slips kept literal: no printed word, symbol, hypothesis or number of the counted statement or of its proof was altered. Layout and class: the article's own class is \texttt{elsarticle} with packages including lineno, mathrsfs, amsfonts, enumerate, latexsym, xy, verbatim, color, hyperref, amsmath, amssymb, xspace, amsthm; the wrapper is the lane's pinned \texttt{amsart} editorial wrapper at 10pt on A4, which restates the theorem-like environments the retained text needs and adds the article's own macro definitions that the retained text needs (\texttt{\textbackslash Bil}, \texttt{\textbackslash Inv}, \texttt{\textbackslash bC}, \texttt{\textbackslash bZ}, \texttt{\textbackslash dg}, \texttt{\textbackslash fL}, \texttt{\textbackslash fg}, \texttt{\textbackslash fl}), with extra wrapper packages (bm, bbm, dsfont, xspace, cancel, mathtools). The delivered numbering is the unit's own. Assets: no retained span contains a figure environment, a graphics-inclusion command, a PDF-page inclusion, a table or a TikZ picture; no image file is copied into this package, the package declares no asset dependency and no image file is redistributed with it. Licence: version arXiv:2507.00349v1 of this article is distributed under the Creative Commons Attribution 4.0 International licence, as recorded in the arXiv licence metadata for this exact version and shown on the abstract page https://arxiv.org/abs/2507.00349v1. A full rights notice is delivered at licenses/arxiv-2507.00349-CC-BY-4.0.txt. Characters: every retained excerpt is pure ASCII, so the delivered engine, XeLaTeX with its default Latin Modern text font, reports no missing character.
\begin{lemm} \label{lemm-6} Let  $\alpha\in H, i\in \bZ$.
	\begin{itemize}
		\item[(a).]
 We have \begin{equation}\label{d-g-4}\alpha(\fl_{i},xt^{-i})=\alpha\left(\fl_2,\Big((\binom i2-\binom {i+1}3)d+\binom i2\Big)xt^{-2}\right).\end{equation}
 	\item[(b).]
$ \alpha(\fl_2,(d^2+d)(x)t^{-2})=0,\forall x\in \dg.$
 	\item[(c).]   $\alpha(\fl_{i},xt^{-i})=0$ if $xt^{-i}\in \fL_{(\mu)}$ with $\mu\ne0, -1$.
 		\item[(d).] $\alpha(\fL_a,\fL_b)=0,$ if $ a+b\ne 0$.
 		\end{itemize}\end{lemm}

\begin{equation}\label{g-g-1}\alpha((a+k(d-1))xt^a,yt^{-a})=\alpha(xt^{a-k},(a-kd)yt^{k-a})+\alpha(\fl_{k},[x,y]t^{-k}).\end{equation}

\begin{equation}\label{o-m}\begin{aligned}\alpha((&a+2k(d-1))xt^a,(a-k-kd)(a-kd)yt^{-a})\\& -\alpha((a+k(d-2))(a+k(d-1))xt^a,(a-2kd)yt^{-a})\\
=&-\alpha(\fl_{k},[x,(a-2kd)(a-kd)y]t^{-k}) -\alpha(\fl_{k},[(a+k(d-2))x,(a-2kd)y]t^{-k})
\\&+\alpha(\fl_{2k},[x,(a-k-kd)(a-kd)y]t^{-2k}).\end{aligned}\end{equation}

\begin{equation}\label{lhd}\begin{aligned} \alpha(2(d-1)&xt^a, d(d+1)yt^{-a})-\alpha((d-2)(d-1)xt^a,-2dyt^{-a})\\
&=2(\alpha((d-1)xt^a, (d+1)dyt^{-a})+\alpha((d-2)(d-1)xt^a,dyt^{-a}))\\
&=2\Big(\alpha((d-1)xt^a, (d-\frac 12)dyt^{-a})+\alpha((d-\frac 12)(d-1)xt^a,dyt^{-a})\Big).\end{aligned}\end{equation}

\begin{equation}\label{Eq2.13}\begin{aligned}&(\pi_0B)(\fl_i,\fl_j)=\delta_{i+j,0}\frac{i^3-i}{12} B(\partial,\partial),\,\,(\pi_0 B)(\fl_i,xt^a)=\delta_{i+a,0}(i^2-i) B(\partial,x),\\ &(\pi_0B)(xt^a,yt^b)=\delta_{a+b,0}\big(aB(x,y)+B(\partial,[x,y])\big),\forall i,j\in \bZ, xt^a,yt^b\in \fg;\end{aligned}\end{equation}

 \begin{lemm}\label{lemm-10}The linear map  $\pi_0$  defined in (\ref{Eq2.13}) is surjective.\end{lemm}% For any $\alpha_0\in H_{(0)}$, we have $\alpha_{0}(\fl_i,dxt^{-i})=0, \alpha_{0}(\fl_i,[x,dy])=0,\alpha_{0}(xt^a,yt^{-a})=-\frac a2\alpha_{0}(\fl_2,[x,y]t^{-2})...

 \begin{proof} Take $\alpha_{0}\in H_{(0)}$. For any $xt^a\in \dg_{(\mu)},yt^{b}\in \dg_{(\lambda)}
 	$  where $a, b\in\frac 1n\bZ, \lambda,\mu\in\bC$, we have
 	$$\alpha_{0}(xt^a,yt^b)=0\text{ if }a+b\ne0\text{ or }\lambda+\mu\ne 0.$$
 	So we take  $xt^a\in \dg_{(\mu)},yt^{-a}\in \dg_{(-\mu)}$. Formulas we will get   actually hold for any $x, y$ even they are not generalized eigenvectors with respect to $d$. Since
 	$(d+1)\dg_{(0)}=\dg_{(0)}$,
 	using Lemma \ref{lemm-6}(b) we know   that $\alpha_{0}(\fl_2  ,d\dg_{(0)}t^{-2})=0$. Furthermore
 	\begin{equation}\label{d}
 		\alpha_{0}(\fl_2  ,d\dg t^{-2})=0; \alpha_0(\fl_2, [x,f(d)y]t^{-2})=\alpha_0(\fl_2,[f(-d)x,y]t^{-2}),\forall f(t)\in \bC[t].\end{equation}
 	Applying this to  (\ref{o-m}), we have
 \begin{equation}\label{2.19}\begin{aligned}\alpha_{0}&\big((a+2k(d-1))xt^a,(a-k-kd)(a-kd)yt^{-a}\big)\\ &-\alpha_{0}\big((a+k(d-2))(a+k(d-1))xt^a,(a-2kd)yt^{-a}\big)\\=&-\binom k2\alpha_{0}\big(\fl_{2},[x,(a-2kd)(a-kd)y]t^{-2}\big)\\
 		&-\binom k2\alpha_{0}\big(\fl_{2},[(a+k(d-2))x,(a-2kd)y]t^{-2}\big)\\&
+\binom {2k}2\alpha_{0}(\fl_{2},[x,(a-k-kd)(a-kd)y]t^{-2})\\
=&-\binom k2\alpha_{0}(\fl_{2},[(a+2kd)(a+kd)x,y]t^{-2})\\
 		&-\binom k2\alpha_{0}(\fl_{2},[(a+2kd)(a+k(d-2))x,y]t^{-2})\\&
+\binom {2k}2\alpha_{0}(\fl_{2},[(a-k+kd)(a+kd)x,y]t^{-2})\\
=&k^2\alpha_{0}(\fl_{2},[a(d+a)x,y]t^{-2})+k^3\alpha_{0}(\fl_{2},[(d+a)(d-1)x,y]t^{-2}).\end{aligned}\end{equation}

%The coefficient of $k^4$ gives $\alpha_{0}(\fl_2, (-[(d-2)x,dy]+[x,d^2y]-2[x,d(d+1)y])t^{-2})=0$.



From (\ref{lhd}), the coefficients of $k^3$ in (\ref{2.19}) give
$$\alpha((d-1)xt^a, (d-\frac 12)dyt^{-a})+\alpha((d-\frac 12)(d-1)xt^a,dyt^{-a})\\
=\frac 12 \alpha_{0}(\fl_2,[(d+a)(d-1)x,y]t^{-2}).
$$

Therefore, $$\alpha(xt^a, (d-\frac 12)dyt^{-a})+\alpha((d-\frac 12)xt^a,dyt^{-a})\\
=\frac 12\alpha_{0}(\fl_2,[(d+a)x,y]t^{-2}).
$$


 Let $B(x,y)=\alpha_{0}(xt^a,dyt^{-a})+\frac 12\alpha_{0}(\fl_2,[(d+a)x,y]t^{-2})$. Then
 $$\begin{aligned}B((d-\frac 12)x,y)&+B(x,(d-\frac 12)y)=\alpha_{0}((d-\frac 12)xt^a,dyt^{-a})+\frac 12\alpha_{0}(\fl_2,[(d+a)(d-\frac 12)x,y]t^{-2})\\ &+\alpha_{0}(xt^a,(d-\frac 12)dyt^{-a})+\frac 12\alpha_{0}(\fl_2,[(d+a)x,(d-\frac 12)y]t^{-2})=0.\end{aligned}$$ Again we have $B(x,y)=0$, that is

 \begin{equation}\label{02.23}\alpha_{0}(xt^a,dyt^{-a})=-\frac 12\alpha_{0}(\fl_2,[(d+a)x,y]t^{-2}).\end{equation}

 Exchanging $x$ and $y$, we deduce that $$\alpha_{0}(dxt^a,yt^{-a})=-\frac 12\alpha_{0}(\fl_2,[x,(d-a)y]t^{-2}).$$

Combining with (\ref{g-g-1}), we deduce that $$\begin{aligned}(a-k)\alpha_0&(xt^a,yt^{-a})-a\alpha_0(xt^{a-k},yt^{k-a})\\
	=&-k\alpha_0(dxt^a,yt^{-a})-k\alpha_0(xt^{a-k},dyt^{k-a})+\alpha_0(\fl_{k},[x,y]t^{-k})\\ =&(-k\frac a2+k\frac{a-k}2+\binom k2)\alpha(\fl_2,[x,y]t^{-2})=-\frac k2\alpha(\fl_2,[x,y]t^{-2}), \forall k\in \bZ,\end{aligned}$$ i.e.,

\begin{equation}\label{02.24}\aligned (a-k)(\alpha_0&(xt^a,yt^{-a})-\frac 12\alpha(\fl_2,[x,y]t^{-2}))\\
 =&a\Big(\alpha_0(xt^{a-k},yt^{k-a})-\frac 12\alpha(\fl_2,[x,y]t^{-2})\Big),\forall xt^a,yt^{-a}\in \fg, k\in \bZ.\endaligned\end{equation}

Now we can define the bilinear form on $B\in (\Bil(\ddot{\fg}))^{\sigma}$ by $$\begin{aligned} B(x,y):=&\frac 1a\Big(\alpha_0(xt^a,yt^{-a})-\frac 12\alpha(\fl_2,[x,y]t^{-2})\Big),\forall xt^a,yt^{-a}\in \fg,a\ne 0; \\ B(\partial,x)=&B(x,\partial)=\frac 12\alpha_0(\fl_2, xt^{-2}),\forall xt^{-2}\in \fg.\end {aligned}$$ It is straightforward to check that $B\in (\Inv(\ddot\fg))^{\sigma}$ and $\alpha_0=\pi_{0}(B)$. Hence $\pi_0$ is surjective.




%$\alpha_{-1}(xt^a,d(d-1)y^{-a})=\alpha((d^2-d)xt^a,yt^a),\forall a\ne 0$.


\end{proof}

\end{document}
